% ============================================================
% PREMIÈRE PARTIE : SEMAINES 1 À 10
% ============================================================

\begin{figure}[H]
\centering

\begin{ganttchart}[
    vgrid,
    hgrid,
    x unit=1cm,
    y unit chart=0.5cm,
    bar height=0.6,
    bar/.style={fill=blue!45},
    bar label font=\footnotesize,
    milestone/.style={fill=red!70},
    milestone label font=\footnotesize,
    title height=1,
    title label font=\bfseries\footnotesize,
    title/.style={fill=gray!25},
    group label font=\bfseries\footnotesize
]{1}{10}

\gantttitle{Semaines 1 à 10}{10} \\
\gantttitlelist{1,...,10}{1} \\

\ganttbar{\tiny Étude bibliographique}{1}{4} \\
\ganttbar{\tiny Rédaction du rapport bibliographique}{1}{4} \\
\ganttbar{\tiny Compréhension de la physique du transfert radiatif}{1}{8} \\
\ganttbar{\tiny Prise en main du code hérité \texttt{embree2TRUST}}{3}{6} \\

\ganttbar{\tiny Développement : correction de la physique}{5}{7} \\
\ganttbar{\tiny Développement : extension 2D}{8}{10} \\
\ganttbar{\tiny Développement : faces réflectives}{10}{10} \\

\ganttbar{\tiny Vérification : numérique}{4}{10} \\
\ganttbar{\tiny Vérification : solutions analytiques}{8}{10} \\

\ganttmilestone{Journée au LJLL}{2} \\
\ganttmilestone{Formation \trust}{3} \\

\end{ganttchart}

\caption{Diagramme de Gantt du déroulement du stage -- semaines 1 à 10.}
\label{fig:gantt_1_10}
\end{figure}
Une journée de séminaire au LJLL a notamment été suivie durant la deuxième semaine du stage
\footnote{\url{https://www.ljll.fr/wp-content/uploads/2026/03/programme_MathNeut_2026.pdf}}.

% ============================================================
% DEUXIÈME PARTIE : SEMAINES 11 À 20
% ============================================================

\begin{figure}[H]
\centering

\begin{ganttchart}[
    vgrid,
    hgrid,
    x unit=1cm,
    y unit chart=0.5cm,
    bar height=0.6,
    bar/.style={fill=blue!45},
    bar label font=\footnotesize,
    milestone/.style={fill=red!70},
    milestone label font=\footnotesize,
    title height=1,
    title label font=\bfseries\footnotesize,
    title/.style={fill=gray!25},
    group label font=\bfseries\footnotesize
]{1}{10}

\gantttitle{Semaines 11 à 20}{10} \\
\gantttitlelist{11,...,20}{1} \\

% S11-13
\ganttbar{\tiny Développement : faces réflectives}{1}{3} \\

% S13-16
\ganttbar{\tiny Développement : corrections numériques}{3}{6} \\

% S11-15
\ganttbar{\tiny Vérification : solutions analytiques}{1}{5} \\

% S13-16
\ganttbar{\tiny Vérification : PyMCRad}{3}{6} \\

% S17
\ganttbar{\tiny Génération des scripts de maillage : \textsc{Salome}}{6}{7} \\

% S18
\ganttbar{\tiny Génération des scripts de simulation : \trust}{8}{8} \\

% S19-20
\ganttbar{\tiny Balayage des simulations : $N_z$, 3 niveaux de puissance}{9}{10} \\
\ganttbar{\tiny Rédaction du rapport}{7}{10} \\

\ganttmilestone{Bonne compréhension du code}{2} \\
\ganttmilestone{Bonne vérification du code}{4} \\
\ganttmilestone{Présentation croisée -- Cadarache}{6} \\

\end{ganttchart}

\caption{Diagramme de Gantt du déroulement du stage -- semaines 11 à 20.}
\label{fig:gantt_11_20}
\end{figure}
